\begin{problem}
    Una persona preparará un postre de frutilla. La receta que tiene indica que por cada \SI{500}{g} de frutilla deberá usar \SI{200}{g} de azúcar. Para elaborar el postre utilizará \SI{100}{g} más de frutilla de lo indicado en la receta. ¿Cuántos gramos de azúcar en total deberá utilizar la persona para mantener la proporción indicada en la receta?

    \begin{enumerate}[label=\Alph*.]
        \item \SI{100}{g}
        \item \SI{140}{g}
        \item \SI{240}{g}
        \item \SI{300}{g}
    \end{enumerate}
\end{problem}
